\section{Conclusion}
% We investigated the optimal language choice for reasoning skeletons in multilingual settings through the Language-Aware Skeleton Exploration Framework (\Lasef). Our extensive experiments empirically demonstrate that while English skeletons are generally effective, they are not a one-size-fits-all solution. We additionally identified a ``Quality–Performance Inversion,'' where linguistically aligned skeletons consistently outperform higher-quality English skeletons for specific low-resource languages. These results highlight that skeleton language is a crucial design variable and suggest that adaptive language exploration is essential for maximizing multilingual reasoning performance.

We analyzed the problem of skeleton language selection in multilingual reasoning through the Language-Aware Skeleton Exploration Framework (\Lasef). Our experiments show that English skeletons are generally effective, but they are not a one-size-fits-all solution. By combining greedy decoding, multi-rollout evaluation, and translation ablation, we further find that skeleton-language effects can be categorized into three patterns: stable cross-lingual, evaluation-dependent, and asymmetric. These results suggest that skeleton language is a context-dependent design variable that cannot be fully explained by generation quality alone, and that effective language combinations should be identified through multi-level exploration.